\documentclass[11pt]{article}
\usepackage[margin=1in]{geometry}
\usepackage{amsmath,amssymb}
\usepackage{enumitem}
\newcommand{\checkmk}{\ensuremath{\checkmark}}

\title{ME18: Kinetic Theory \& Calorimetry --- Walkthrough}
\author{}\date{}

\begin{document}\maketitle

\section*{Core equations}
\begin{itemize}[leftmargin=*]
\item $PV = nRT = Nk_BT$, $R = 8.314$, $N_A = 6.022\times10^{23}$, $k_B = R/N_A$
\item Combined: $P_1V_1/T_1 = P_2V_2/T_2$ (with $n_1,n_2$ if mass changes)
\item $\langle KE\rangle = \tfrac32 k_BT$; total $K_{\text{tr}} = \tfrac32 nRT$
\item Equipartition: $U = \tfrac{f}{2}nRT$
\item $v_{\text{avg}} = \sqrt{8RT/(\pi M)}$, $v_{\text{rms}} = \sqrt{3RT/M}$, $v_p = \sqrt{2RT/M}$ (M in kg/mol)
\end{itemize}

\section*{Q1 --- Volume at STP}
\emph{Direct application of $PV=nRT$ at 1 atm and 273 K.} Expect $\sim$22.4 L/mol scaled by 0.8.

$V = nRT/P = 0.8(8.314)(273)/101325 = 0.01792\ \text{m}^3 = \boxed{17.92\ \text{L}}$ \checkmk

\section*{Q2 --- Combined gas law, no mass change}
\emph{Sealed sample: $P_2=1.9P_1$, $T_2=1.85T_1$.} With mass fixed, $PV/T$ is conserved.

$V_2 = V_1(P_1/P_2)(T_2/T_1) = 1\cdot(1/1.9)(1.85) = \boxed{0.9737\ \text{m}^3}$ \checkmk

\section*{Q3 --- Solve for new temperature}
\emph{Sealed sample shrinks both in pressure ($\times 0.9$) and volume ($\times 0.7$); temperature must drop accordingly.}

$T_2 = T_1(P_2V_2/(P_1V_1)) = 297(0.9)(0.7) = \boxed{187.11\ \text{K}}$ \checkmk

\section*{Q4 --- Tire with leak \& warming}
\emph{Tire warms 19 $^\circ$C while losing 9\% of its gas.} Convert gauge $\to$ absolute pressure first; mole count changes, so use the generalized form with $n_2 = 0.91 n_1$ and constant $V$.

$P_1 = 311{,}466$ Pa abs, $T_1 = 278.15$ K, $T_2 = 297.15$, $n_2 = 0.91n_1$.
$P_2 = P_1(0.91)(297.15/278.15) = 302810$ Pa abs.
$P_{\text{gauge}} = 202810/6895 = \boxed{29.41\ \text{psi}}$ \checkmk

\section*{Q5 --- Molecules in 77.1 ag of C$_4$H$_8$OH}
\emph{Three-step conversion: attograms $\to$ g $\to$ mol (via molar mass) $\to$ molecules (via $N_A$).} Molecular formula C$_4$H$_8$OH = C$_4$H$_9$O.

$M_{\text{C}_4\text{H}_9\text{O}} = 73.116$ g/mol; $n = 7.71\times10^{-17}/73.116 = 1.0545\times10^{-18}$ mol; $N = nN_A = \boxed{6.351\times10^5}$ \checkmk

\section*{Q6 --- Moles of AlCl$_3$ in 501 g}
\emph{Compute molar mass from atomic weights, then $n=m/M$.}

$M_{\text{AlCl}_3} = 133.33$ g/mol; $n = 501/133.33 = \boxed{3.757\ \text{mol}}$ \checkmk

\section*{Q7 --- Temperature from translational KE}
\emph{Translational KE of $n$ moles is $\tfrac32 nRT$ regardless of molecular mass.} Solve for $T$.

$T = 2K/(3nR) = 2(64356)/(3\cdot 8\cdot 8.314) = \boxed{645.02\ \text{K}}$ \checkmk

\section*{Q8 --- Total KE of 7.6 mol at 90 $^\circ$C}
\emph{Plug into $K=\tfrac32 nRT$; convert $^\circ$C $\to$ K first.}

$K = 1.5(7.6)(8.314)(363.15) = \boxed{34{,}407\ \text{J}}$ \checkmk

\section*{Q9 --- Average speed of O$_2$}
\emph{Maxwell-Boltzmann mean speed $\sqrt{8RT/(\pi M)}$, with $M$ for \emph{molecular} O$_2$ = 0.032 kg/mol.}

$v_{\text{avg}} = \sqrt{8(8.314)(355.05)/(\pi\cdot 0.032)} = \boxed{484.67\ \text{m/s}}$ \checkmk

\section*{Q10 --- Average $v_x$}
\emph{Speeds are nonzero, but a signed velocity component averages to zero by symmetry --- for every molecule moving $+x$, another moves $-x$.}

$\langle v_x\rangle = \boxed{0}$ (symmetric distribution).

\section*{Q11 --- RMS speed at 39 $^\circ$C}
\emph{$v_{\text{rms}} = \sqrt{3RT/M}$ --- the speed appearing in $\tfrac12 m v_{\text{rms}}^2 = \tfrac32 k_BT$.}

$v_{\text{rms}} = \sqrt{3(8.314)(312.15)/0.032} = \boxed{493.15\ \text{m/s}}$ \checkmk

\section*{Q12 --- Internal energy of helium}
\emph{Helium monatomic: only 3 translational dof, so $U = \tfrac{3}{2}nRT$. Convert $^\circ$C $\to$ K.}

He monatomic, $f=3$. $U = 1.5(2.4)(8.314)(428.15) = \boxed{12.81\ \text{kJ}}$ \checkmk

\section*{Q13 --- Internal energy change with $f=2$}
\emph{Equipartition: $\tfrac12 RT$ per mole per degree of freedom; with $f=2$, $\Delta U = nR\Delta T$.}

$\Delta U = (f/2)nR\Delta T = 1\cdot 6.2(8.314)(53) = \boxed{2.732\ \text{kJ}}$ \checkmk

\section*{Q14 --- Definition of $k_B$}
\emph{Bridges per-molecule and per-mole forms: $nR = Nk_B \Rightarrow k_B = R/N_A$.}

$k_B = R/N_A$ --- ``the gas constant \textbf{divided by} Avogadro's number.''

\section*{Q15 --- Same $Q$ into monatomic vs diatomic gas}
\emph{At constant volume, $Q=\Delta U = (f/2)nR\Delta T$. Diatomic has $f=5$ vs monatomic $f=3$, so the diatomic gas absorbs $Q$ into more ``energy buckets'' and warms less.} For $v_{\text{rms}}\propto\sqrt{T/M}$ we need molecular masses, which aren't given.

$\Delta T = 2Q/(fnR)$ with $f_A=3 < f_B=5$ $\Rightarrow$ A's $T$ \textbf{higher}. $v_{\text{rms}}\propto\sqrt{T/M}$ but $M$ unknown $\Rightarrow$ \textbf{impossible to determine}.

\section*{Q16 --- Ordering of speed measures}
\emph{For Maxwell-Boltzmann, $v_{\text{rms}}^2 = \langle v^2\rangle$ weights fast molecules more, while $v_p$ (peak of distribution) sits below the mean because of the long high-speed tail.}

Order: $v_p < v_{\text{avg}} < v_{\text{rms}}$. Highest = \textbf{rms}; lowest = \textbf{most probable}.

\bigskip\noindent\textbf{All numeric answers match source key.}
\end{document}
